\documentclass[UTF8,a4paper,12pt]{ctexart}
\usepackage[margin=2cm]{geometry}
\usepackage{amsmath,amssymb,fancyhdr}
\pagestyle{fancy}\fancyhf{}\fancyfoot[C]{\thepage}\renewcommand{\headrulewidth}{0pt}
\setlength{\parindent}{0pt}\setlength{\parskip}{5pt}
\newcommand{\src}[1]{{\small #1}\par}
\begin{document}
\begin{center}{\LARGE 高一数学竞赛选题}\par 集合与基本不等式·开学三周拓展\par 时间：90分钟\qquad 满分：100分\end{center}
姓名：\underline{\hspace{3cm}}\qquad 班级：\underline{\hspace{3cm}}

\textbf{1．（16分）}\src{改编自2011年全国高中数学联合竞赛一试A卷第1题}
集合 $A=\{a_1,a_2,a_3,a_4\}$ 含有四个不同的实数。将 $A$ 的每个三元素子集中的元素相加，所得各个和构成集合 $B=\{-1,3,5,8\}$。求集合 $A$，并说明理由。
\par\vspace*{2.4cm}

\textbf{2．（18分）}\src{改编自2014年全国高中数学联赛一试第2题}
设
\[S=\left\{\frac3a+b\ \middle|\ a,b\in\mathbb R,\ 1\le a\le b\le2\right\}.\]
求集合 $S$ 中的最大值 $M$、最小值 $m$ 及 $M-m$，并证明所求结论。
\par\vspace*{2.4cm}

\textbf{3．（20分）}\src{改编自2011年全国高中数学联合竞赛一试A卷第3题}
正实数 $a,b$ 满足
\[\frac1a+\frac1b\le2\sqrt2,\qquad (a-b)^2=4(ab)^3.\]
求所有满足条件的有序数对 $(a,b)$，并写出推导过程。
\par\vspace*{2.4cm}
\newpage
\textbf{4．（22分）}\src{选自1972年第14届国际数学奥林匹克（IMO）第1题（中文转述）}
集合 $A$ 由十个不同的两位正整数构成。证明：一定存在 $A$ 的两个非空子集 $B,C$，使得 $B\cap C=\varnothing$，并且 $B$ 中所有元素的和等于 $C$ 中所有元素的和。
\par\vspace*{6cm}

\textbf{5．（24分）}\src{摘编自2008年第49届国际数学奥林匹克（IMO）第2题（a）}
实数 $x,y,z$ 均不等于 $1$，且满足 $xyz=1$。证明：
\[\frac{x^2}{(x-1)^2}+\frac{y^2}{(y-1)^2}+\frac{z^2}{(z-1)^2}\ge1.\]
\par\vspace*{6cm}
\end{document}
